% answers_short.tex —— 2005 年·普通高等学校招生考试·重庆卷·理科 · 简短答案速查
% 由 tools/make_short_answers.py 自 answers.tex 的 \gksolhead/\gkanswer 提取，
% 并与 2005/2005_数学真题_重庆理科_答案要点.md 逐题比对后生成，勿手改（改源 MD 后重跑）。

\begin{gknotice}
\textbf{使用说明}：本节只给最终答案便于快速核对，完整推导见解析版。答案为编写者独立推导，\textbf{非官方标准答案}。
\end{gknotice}

\gkpart{选择题}
\begin{center}\begin{tabular}{ccccccccccc}
\toprule
题号 & 1 & 2 & 3 & 4 & 5 & 6 & 7 & 8 & 9 & 10 \\
\midrule
答案 & A & A & D & C & C & B & B & B & A & C \\
\bottomrule
\end{tabular}\end{center}

\gkpart{填空题}
\begin{itemize}[leftmargin=2.6em]
  \item[\textbf{11.}] $(0, 3)$（或 $\{x \mid 0 < x < 3\}$）
  \item[\textbf{12.}] $\pm 1$（或 $1$ 或 $-1$）
  \item[\textbf{13.}] $1$
  \item[\textbf{14.}] $-3$
  \item[\textbf{15.}] $\dfrac{45}{128}$
  \item[\textbf{16.}] ②③⑤
\end{itemize}

\gkpart{解答题}
\begin{itemize}[leftmargin=2.6em]
  \item[\textbf{17.}] $a = \pm\sqrt{15}$。
  \item[\textbf{18.}] (1) $\dfrac{2}{3}$；(2) $\xi$ 的分布列为 $P(\xi=0)=\dfrac{1}{3}, P(\xi=10)=\dfrac{2}{5}, P(\xi=20)=\dfrac{1}{15}, P(\xi=50)=\dfrac{2}{15}, P(\xi=60)=\dfrac{1}{15}$；数学期望 $\mathrm{E}\xi = 16$（元）
  \item[\textbf{19.}] 当 $0 \leqslant a \leqslant 4$ 时，函数 $f(x)$ 的极值点个数为 0；当 $a < 0$ 或 $a > 4$ 时，函数 $f(x)$ 的极值点个数为 2。
  \item[\textbf{20.}] (1) $1$；(2) $\dfrac{\sqrt{2}}{2}$
  \item[\textbf{21.}] (1) $\dfrac{x^2}{3} - y^2 = 1$；(2) $\left(-1, -\dfrac{\sqrt{195}}{15}\right) \cup \left(-\dfrac{\sqrt{3}}{3}, -\dfrac{1}{2}\right) \cup \left(\dfrac{1}{2}, \dfrac{\sqrt{3}}{3}\right) \cup \left(\dfrac{\sqrt{195}}{15}, 1\right)$
  \item[\textbf{22.}] (1) 见解析版；(2) 见解析版
\end{itemize}
